\documentclass[10pt,a4paper]{amsart}
\usepackage[margin=19mm]{geometry}
\usepackage{amsmath,amssymb,mathtools}
\usepackage[scr=rsfs,cal=euler]{mathalfa}
\usepackage{xfrac}
\usepackage[unicode,hidelinks,bookmarks=false]{hyperref}
\newcommand{\ie}{i.e.\ }
\newcommand{\cf}{cf.\ }
\newcommand{\resp}{respectively\ }
\newcommand{\Subsection}{\subsection}
\providecommand{\nobreakdash}{\nobreak\hbox{-}}
\setlength{\emergencystretch}{2em}
\multlinegap0pt
\usepackage{bm}
\usepackage{bbm}
\usepackage{dsfont}
\usepackage{xspace}
\usepackage{cancel}
\usepackage{mathtools}
\usepackage{cleveref}
\usepackage{amsthm}
\makeatletter
\@ifundefined{theorem}{\newtheorem{theorem}{Theorem}}{}
\@ifundefined{lemma}{\newtheorem{lemma}[theorem]{Lemma}}{}
\makeatother
\newcommand{\Ex}{\mathbb{E}}
\newcommand{\cF}{\mathcal{F}}
\newcommand{\cP}{\mathcal{P}}
\title[Pancyclicity of highly connected graphs]{Pancyclicity of highly connected graphs}
\author{Shoham Letzter}
\date{Source: arXiv:2306.12579v3 (CC BY 4.0)}
\begin{document}
\maketitle

\noindent\emph{Source and licence.} Pancyclicity of highly connected graphs. Version arXiv:2306.12579v3 (CC BY 4.0), distributed under the Creative Commons Attribution 4.0 International licence, as recorded in the arXiv licence metadata for this exact version and shown on the abstract page https://arxiv.org/abs/2306.12579v3. Full rights notice: licenses/arxiv-2306.12579-CC-BY-4.0.txt.\par\smallskip

\noindent\textit{Editorial scope.} Counted unit: the Lemma whose source label is \texttt{lem:odd-cycles} in the article (source0.tex lines 1026--1030, source label \texttt{lem:odd-cycles}), delivered together with the article's own complete original proof of it (source0.tex lines 1078--1119, one \texttt{proof} environment of 42 lines). No mathematics is written, repaired or re-derived: statement and proof are the article's own text line for line. The proof is detached from the statement in the source: the article states the result at lines 1026--1030 and prints its proof at lines 1078--1119, and the proof environment's own header names the statement, so the association is the article's own. Required context retained uncounted: none. Every definition, display and label of those lines is retained unchanged. Omitted from this package and not redistributed: 1--1025, 1031--1077, 1120--1306. Nothing inside a retained excerpt is omitted or altered: every retained body is delivered exactly as the article's lines are written, so no omission receipt of any kind accompanies this package. Citations: the retained excerpts carry no citation command at all, so no bibliography entry of the article is transcribed and no reference is redistributed. Reference adaptation: none was needed and none was made; every reference inside the delivered unit resolves inside it, so no unresolved reference or citation is produced. Printed slips kept literal: no printed word, symbol, hypothesis or number of the counted statement or of its proof was altered. Layout and class: the article's own class is \texttt{article} with packages including geometry, amsthm, amsmath, amssymb, mathtools, graphicx, hyperref, color, thm-restate, enumitem, subcaption, caption, natbib, setspace; the wrapper is the lane's pinned \texttt{amsart} editorial wrapper at 10pt on A4, which restates the theorem-like environments the retained text needs and adds the article's own macro definitions that the retained text needs (\texttt{\textbackslash Ex}, \texttt{\textbackslash cF}, \texttt{\textbackslash cP}), with extra wrapper packages (bm, bbm, dsfont, xspace, cancel, mathtools, cleveref). The delivered numbering is the unit's own. Assets: no retained span contains a figure environment, a graphics-inclusion command, a PDF-page inclusion, a table or a TikZ picture; no image file is copied into this package, the package declares no asset dependency and no image file is redistributed with it. Licence: version arXiv:2306.12579v3 of this article is distributed under the Creative Commons Attribution 4.0 International licence, as recorded in the arXiv licence metadata for this exact version and shown on the abstract page https://arxiv.org/abs/2306.12579v3. A full rights notice is delivered at licenses/arxiv-2306.12579-CC-BY-4.0.txt. Characters: every retained excerpt is pure ASCII, so the delivered engine, XeLaTeX with its default Latin Modern text font, reports no missing character.
		\begin{lemma} \label{lem:odd-cycles}
			Let $0 < \delta \ll 1$ and let $\alpha, n \ge 1$ be integers.
			Suppose that $G$ is an $n$-vertex graph with $\delta(G) > \alpha(G) = \alpha$.
			Then there is a subgraph $H$ of $G$ with at least $\delta n \alpha$ edges, such that each edge $e$ in $H$ is contained in a cycle $C$ of length $3$ or $5$ satisfying $V(C) \cap V(H) = V(e)$.
		\end{lemma}

		\begin{proof}[Proof of \Cref{lem:odd-cycles}]
			Let $\eta$ be a constant satisfying $\delta \ll \eta \ll 1$.
			For each vertex $u$ let $M(u)$ be a maximum matching in $N(u)$, and let $I(u) := N(u) - M(u)$ (so $I(u)$ is independent).
			We define sets $U_1, \ldots, U_4$ as follows. 
			% unloaded sounds weird
			\begin{itemize}
				\item
					Let $U_1$ be the set of vertices $u$ satisfying $|M(u)| \ge \alpha/12$.
				\item
					Let $U_2$ be the set of vertices $u \notin U_1$ for which there exists a triangle $(uvw)$, where $v \notin U_1$. 
				\item
					Let $U_3$ be the set of vertices $u \notin U_1 \cup U_2$ for which there exist three vertices $v,w,x \neq u$ such that $(vwx)$ is a triangle, $v \notin U_1$, and $uw$ is an edge.
				\item
					Let $U_4 = V(G) - (U_1 \cup U_2 \cup U_3)$.
			\end{itemize}
			We show that $|U_4| \le n/3$.
			Indeed, for each $u \in U_4$ let $T_u$ be a triangle that contains $u$ (such a triangle exists by $\delta(G) > \alpha$). 
			We claim that the triangles $T_u$, with $u \in U_4$, are pairwise vertex-disjoint.
			To see this, consider distinct $u_1,u_2 \in U_4$, and write $T_u = (u_1v_1w_1)$ and $T_{u_2} = (u_2v_2w_2)$. Notice that $v_1,w_1,v_2,w_2 \in U_1$, because $u_1,u_2 \notin U_2$. Thus, if $T_{u_1}$ and $T_{u_2}$ share a vertex, we may assume $v_1 = v_2$. This is a contradiction to $u \notin U_4$, as $(u_2v_2w_2)$ is a triangle, $u_2 \notin U_1$, and $u_1v_2$ is an edge.
			This shows that the triangles $T_u$, with $u \in U_4$, are pairwise vertex-disjoint, which implies the desired inequality $|U_4| \le n/3$.

			Let $\cF$ be a family of ordered triples and quintuples, defined as follows. 
			It will be useful to note that for distinct $u,v \notin U_1$ either $|I(u) \cap I(v)| \ge \alpha/12$ or there is a matching of size at least $\alpha/12$ whose edges have one end in $I(u)$ and the other in $I(v)$; this follows from $I(u)$ and $I(v)$ being independent sets of size at least $2\alpha/3$.
			\begin{itemize}
				\item
					For every $u \in U_1$ and $xy \in M(u)$, add the ordered triples $(u,x,y)$ and $(u, y, x)$ to $\cF$.
				\item
					For every $u \in U_2$, fix $v, w$ such that $v \notin U_1$ and $(uvw)$ is a triangle.
					Because $I(u)$ and $I(v)$ are two independent sets of size at least $2\alpha/3$.
					If $|I(u) \cap I(v)| \ge \alpha/12$, add the ordered triple $(u,x,v)$ to $\cF$ for each $x \in I(u) \cap I(v)$. 
					Otherwise there is a matching $M$ of size at least $\alpha/12$ between $I(u)$ and $I(v)$. Add to $\cF$ the ordered quintuple $(u, x, y, v, w)$, for each edge $xy \in M$ with $x \in I(u)$ and $y \in I(v)$.
					% need to remove
				\item
					For every $u \in U_3$, fix $v, w, x \neq u$ such that $(v w x)$ is a triangle, $v \notin U_1$, and $uw$ is an edge. If $|I(u) \cap I(v)| \ge \alpha/12$, add $(u, y, v, x, w)$ with $y \in I(u) \cap I(v)$ to $\cF$. Otherwise, let $M$ be an $I(u)-I(v)$ matching of size at least $\alpha/12$, and add $(u, y, z, v, w)$ to $\cF$, for all $yz \in M$ (where $y \in I(u)$ and $z \in I(v)$). 
			\end{itemize}
			Notice that for every $u \notin U_4$ there are at least $\alpha/12$ ordered tuples in $\cF$ whose first element is $u$, showing that $|\cF| \ge 2n/3 \cdot \alpha/12 = \alpha n/18$. Moreover, for every two vertices $x,y$, there is at most one ordered tuple in $\cF$ whose first element is $x$ and the second $y$.

			Let $X$ be a random set of vertices, obtained by including each vertex with probability $1/2$, independently. 
			Let $\cP$ be the set of ordered pairs $(x,y)$ such that $x,y \in X$ and $\cF$ contains a tuple whose first element is $x$ and the second is $y$, and whose other elements are not in $X$.
			Then $\Ex(|\cP|) \ge |\cF| \cdot 2^{-5} \ge \alpha n / 576$. Fix an outcome of $X$ such that $|\cP| \ge \alpha n / 576$.
			Let $H$ be the graph on vertex set $X$, with edges $xy$ such that at least one of $(x,y)$ and $(y,x)$ is in $\cP$. Then $e(H) \ge |\cP|/2 \ge \alpha n / 1152$. The graph $H$ satisfies the requirements of the lemma.
		\end{proof}

\end{document}
